% =====================================================================
%  25.3  Evert, Link, Cut, and Path Aggregates -- the full LCT.
%  Exports: sec:c23-makeroot  def:c23-makeroot  def:c23-evert-agg
%           thm:c23-path-query  lst:c23-evert  lst:c23-pathsum
%  WORKED example only -- NO solvedproblem (the in-section solved count
%  stays 2; the flagship dynamic-path problem is the chapter Summit).
%  The worked task is a self-contained "dynamic tree vertex add / path sum":
%  link, cut, add-to-vertex, query sum along a path, fully traced on a
%  4-vertex chain so the reader sees makeroot expose a path and the reverse
%  lazy tag do its work. Modelled on the standard rooted-then-everted LCT.
% =====================================================================
\section{Evert, Link, Cut, and Path Aggregates}\label{sec:c23-makeroot}

\begin{previouslybox}[Ch.~22 needed a fixed root; the LCT gets to pick one per query]
In \(\secref{c23-prefer}\) the forest stayed rooted and \code{access} could only
expose a path \emph{downward} from that fixed root. To link two arbitrary trees,
cut a single edge, or query the path between two nodes that are neither an
ancestor of the other, we need one more move: \emph{reroot the represented tree at
any node}. Chapter~22 sidestepped this by choosing a root once in its DFS
(\(\secref{c16-decomp}\)); the LCT does it on demand with \code{makeroot}, and the
price is a single lazy flag. This section adds \code{makeroot}, the reverse tag,
\code{link}/\code{cut}, and a path aggregate -- completing the tool the Summit
will ride on a block-cut tree.
\end{previouslybox}

\begin{topicintro}
\code{makeroot(v)} = \code{access(v)} followed by flipping a \emph{reverse} lazy tag
on \(v\)'s aux tree. Since that aux tree currently stores the depth-ordered
root-to-\(v\) path, reversing it makes the old root the deepest node and \(v\) the
new root -- inverting the whole path's parent relations at once, exactly like
reversing a sequence in an implicit splay. Once you can evert, the three composite
operations are one-liners: \code{link(u,v)} everts \(u\) then hangs it under \(v\);
\code{cut(u,v)} everts \(u\), exposes \(v\), and checks the two are adjacent; a path
aggregate \code{query(u,v)} everts \(u\), accesses \(v\), and reads \(v\)'s splay
fold. The fold is any monoid (sum, min, max, gcd) carried up the aux trees by
\code{pull}, and pushed down through the reverse tag by \code{pushdown}.

\textbf{Prerequisites.} \(\secref{c23-prefer}\) -- \code{access}, \code{splay},
\code{is\_root}, \code{findroot}. \(\secref{c23-splay}\) -- lazy-tag-free splaying,
now augmented. The monoid/segment-fold idea is Ch.~11 (\chref{c14}); nothing about it
is re-taught.

\textbf{After this section you can}
\begin{itemize}[nosep]
  \item write \code{makeroot} as \code{access} + a reverse tag, and explain why
        reversal reroots; \hfill\lv{5}
  \item implement \code{link}/\code{cut}/\code{query} from \code{makeroot}; \hfill\lv{5}
  \item carry a vertex-sum (or min) aggregate through rotations and lazy flips; \hfill\lv{5}
  \item trace a dynamic path-sum: evert, expose, fold, read. \hfill\lv{5}
\end{itemize}
\end{topicintro}

\begin{definition}{Evert and the reverse lazy tag}{c23-makeroot}
\code{makeroot(v)} performs \code{access(v)} (so the root-to-\(v\) path is one aux
splay with \(v\) at its root and depth increasing left-to-right) and then applies a
\term{reverse} tag to \(v\)'s splay: swap its two children and mark the tag to be
pushed down. Inverting the in-order sequence inverts the path's direction, so the
old root becomes a leaf and \(v\) becomes the represented root. The tag is a plain
lazy propagation -- identical to reversing a segment in an implicit splay.
\end{definition}

\begin{lstlisting}[caption={Adding evert power to the LCT core: pull, apply-reverse, pushdown, makeroot, link, cut, and a path-sum query. Children ls/rs, aggregate sum[], per-node val[], lazy rev[]. Call pushdown top-down inside splay before rotating. ASCII only}, label={lst:c23-evert}]
int ch[MAXN][2], fa[MAXN], rev[MAXN];
long long val[MAXN], sum[MAXN];
void pull(int x){ sum[x] = val[x] + sum[ch[x][0]] + sum[ch[x][1]]; }
void applyRev(int x){ if(!x) return; swap(ch[x][0], ch[x][1]); rev[x]^=1; }
void pushdown(int x){ if(rev[x]){ applyRev(ch[x][0]); applyRev(ch[x][1]); rev[x]=0; } }
// splay() from lst:c23-lct-core, but push the whole ancestor chain first:
void pushPath(int x){ if(!isRoot(x)) pushPath(fa[x]); pushdown(x); }
//   ... then in splay(x): call pushPath(x) before each rot(...) group.
void makeroot(int x){ access(x); applyRev(x); }
int findroot(int x){ access(x); while(ch[x][0]){ pushdown(x); x=ch[x][0]; } splay(x); return x; }
void link(int u,int v){ makeroot(u); if(findroot(v)!=u) fa[u]=v; }
void cut(int u,int v){
    makeroot(u); access(v);            // path u->v is v's aux tree, u is leftmost
    if(ch[v][0]==u && !ch[u][1]){ ch[v][0]=0; fa[u]=0; pull(v); }
}
long long query(int u,int v){          // sum of val on the simple path u..v
    makeroot(u); access(v);            // u..v now a single aux tree rooted at v
    return sum[v];
}
void addVal(int u,long long d){ access(u); val[u]+=d; pull(u); }
\end{lstlisting}

\begin{theorem}{Evert-then-access exposes a path, and the fold answers it}{c23-path-query}
After \code{query(u,v)} = \code{makeroot(u)} then \code{access(v)}, the auxiliary
splay tree rooted at \(v\) contains exactly the vertices of the simple path
\(u\leadsto v\) in the represented tree, and \(\mathit{sum}[v]\) equals the sum of
their values.
\end{theorem}
\begin{proofbox}
\code{makeroot(u)} makes \(u\) the represented root (\(\defref{c23-makeroot}\)), so
\(v\)'s ancestors run up to \(u\). \code{access(v)} then makes the current
root-to-\(v\) path -- which is precisely \(u\leadsto v\) -- a single preferred path,
i.e.\ one aux splay with \(v\) splayed to the top (\(\defref{c23-access}\)); no
other vertex is in that splay, because access detaches every right child along the
spine, leaving only the path. The in-order sequence of the splay is the path by
depth, and \code{pull} maintains \(\mathit{sum}[x]=\mathit{val}[x]+
\mathit{sum}[\mathit{ls}]+\mathit{sum}[\mathit{rs}]\) through every rotation and
lazy reversal, so the aggregate of the root \(v\) is the sum over the whole splay,
i.e.\ over the path. \(\Box\)
\end{proofbox}

\begin{complexitybox}[Evert, link, cut, query]
\textbf{Per operation} \(\bigO{\log n}\) amortised -- \code{makeroot} is one
\code{access} plus an \(O(1)\) tag; \code{link}/\code{cut}/\code{query} are
constant many accesses/splays. \textbf{Space} \(\bigO{n}\). The reverse tag costs
nothing asymptotically: it is one swap at a node, propagated lazily, and any
aggregate that is invariant under reordering (sum, min, max, gcd, xor) needs no
order-aware fix-up -- only an order-\emph{sensitive} fold would have to also flip
its combine direction inside \code{applyRev}.
\end{complexitybox}

% ---------------------------------------------------------------
\subsection{Worked Example: dynamic vertex-add / path-sum}
% ---------------------------------------------------------------

\begin{problemstatement}[A self-contained dynamic-tree exercise (worked, not a judge problem)]
A forest of weighted vertices starts as isolated nodes. Support online:
\(\texttt{link}(u,v)\) join two trees by an edge; \(\texttt{cut}(u,v)\) remove an
edge; \(\texttt{add}(u,d)\) add \(d\) to vertex \(u\)'s value; \(\texttt{query}(u,v)\)
print the sum of values along the unique path \(u\leadsto v\). All in
\(O(\log n)\) amortised.
\end{problemstatement}

\begin{insightbox}[The whole program is \code{makeroot} + \code{access}]
Add is \code{access(u)} then bump \(\mathit{val}[u]\). Link, cut and query all start
by \code{makeroot}-ing one endpoint so the path becomes root-to-the-other and a
final \code{access} folds it. Once you internalise ``\code{makeroot(u); access(v)}
turns \(u\leadsto v\) into \(v\)'s splay'', every dynamic path aggregate is the same
two lines -- only the \code{pull} monoid changes (sum here, min at the Summit).
\end{insightbox}

\begin{tracebox}[Build a path \(1\!-\!2\!-\!3\!-\!4\), then query and reroot]
Values \(v_1{=}1,v_2{=}2,v_3{=}3,v_4{=}4\).\\
\(\texttt{link}(1,2)\): \code{makeroot(1)} (\(1\) alone, tag on a singleton is a
no-op after pushdown), \(\mathit{fa}[1]{=}2\).\\
\(\texttt{link}(2,3)\): \code{makeroot(2)}; \(2\)'s path is \(2\!-\!1\) reversed to
root \(2\); hang \(2\) under \(3\): \(\mathit{fa}[2]{=}3\). Tree now \(3\!-\!2\!-\!1\).\\
\(\texttt{link}(3,4)\): \code{makeroot(3)} roots the chain at \(3\); \(\mathit{fa}[3]{=}4\).
Represented chain \(4\!-\!3\!-\!2\!-\!1\).\\
\(\texttt{query}(1,4)\): \code{makeroot(1)} re-roots so \(1\) is root
(\(1\!-\!2\!-\!3\!-\!4\)); \code{access(4)} exposes path \(1\!-\!2\!-\!3\!-\!4\) as
\(4\)'s splay; \(\mathit{sum}[4]=1{+}2{+}3{+}4=\boxed{10}\).\\
\(\texttt{add}(3,5)\): \code{access(3)}; \(\mathit{val}[3]=8\); ancestors' sums fixed.\\
\(\texttt{query}(1,4)\) again: path sum \(1{+}2{+}8{+}4=\boxed{15}\).\\
\(\texttt{cut}(2,3)\): \code{makeroot(2)}; \code{access(3)}; check \(\mathit{ch}[3][0]{=}2\)
and \(2\) has no right child; detach \(\Rightarrow\) trees \(\{1,2\}\) and \(\{3,4\}\).\\
\(\texttt{query}(1,2)\): \(\boxed{3}\) (still connected); \(\texttt{query}(1,4)\) now
returns nothing useful -- guard with \(\mathit{findroot}(1)\ne\mathit{findroot}(4)\).
\end{tracebox}

\begin{lstlisting}[caption={The worked path-sum driver (main only; reuse lst:c23-evert verbatim). Queries mirror YS ``dynamic tree vertex add path sum''. ASCII only}, label={lst:c23-pathsum}]
int main(){
    int n,q;  scanf("%d %d",&n,&q);
    for(int i=1;i<=n;i++){ val[i]=/* read */; sum[i]=val[i]; }
    while(q--){
        char op[8]; scanf("%s",op);
        if(op[0]=='l'){ int u,v; scanf("%d %d",&u,&v); link(u,v); }
        else if(op[0]=='c'){ int u,v; scanf("%d %d",&u,&v); cut(u,v); }
        else if(op[0]=='a'){ int u; long long d; scanf("%d %lld",&u,&d); addVal(u,d); }
        else { int u,v; scanf("%d %d",&u,&v);
               if(findroot(u)!=findroot(v)) printf("disconnected\n");
               else printf("%lld\n", query(u,v)); }
    }
}
\end{lstlisting}

\begin{warningbox}[Push the tags BEFORE you rotate -- every LCT bug starts here]
\code{pushPath(x)} must run \emph{top-down} over the ancestor chain before
\code{splay(x)}'s rotations, or a pending reverse tag sits above \(x\) and the
rotation reads a stale child pointer -- the tree silently loses or duplicates path
segments (wrong answer only on paths that cross an everted region, so it hides for
dozens of operations). Pair it with the mirror discipline in \code{makeroot}: flip
\(\mathit{rev}\) and swap children \emph{together}, and never call \code{pull}
between the swap and a needed \code{pushdown}.
\end{warningbox}

\begin{hardeningbox}
\begin{itemize}[nosep,leftmargin=1.4em]
  \item \textbf{Guard link/cut.} \code{link} should skip when \(u,v\) are already
        connected (\(\mathit{findroot}(v)==u\)) or you create a cycle and corrupt the
        structure; \code{cut} must verify \(\mathit{ch}[v][0]==u\) and \(\mathit{ch}[u][1]==0\)
        before detaching, so a non-edge cut is a harmless no-op.
  \item \textbf{Order-sensitive folds need a second aggregate.} Sum/min/max/gcd are
        commutative, so one \code{sum[]} suffices. For a fold like
        ``string along the path'' you must store both forward and reversed folds and
        swap them in \code{applyRev}.
  \item \textbf{findroot's pushdown.} The left-descent in \code{findroot} has to
        \code{pushdown} at each step (see \(\coderef{c23-evert}\)); a pending reversal
        makes the ``leftmost'' pointer wrong and returns the wrong root.
\end{itemize}
\end{hardeningbox}

\begin{recallbox}[25.3 in one breath]
\emph{\code{makeroot(u)} = \code{access(u)} + a reverse tag, turning \(u\) into the
represented root. Then \code{link}/\code{cut}/\code{query} are ``evert one endpoint,
access the other'': \code{query(u,v)} folds \(v\)'s aux splay, which after
\code{makeroot(u); access(v)} is exactly the path \(u\leadsto v\)
(\(\thmref{c23-path-query}\)). Push lazy tags top-down before rotating or the structure
rots out. This full LCT is what the chapter Summit mounts on a block-cut tree.}
\selftest{
  \item What two calls expose path \(u\leadsto v\) as one aux splay? \emph{(\code{makeroot(u)} then \code{access(v)}, \(\thmref{c23-path-query}\).)}
  \item Why is a reverse tag enough to reroot? \emph{(the aux splay is the depth-ordered path; inverting the sequence inverts parent direction, Def.~\defref{c23-makeroot}.)}
  \item Which aggregates are free to fold, and which need extra care? \emph{(commutative ones -- sum/min/max/gcd/xor -- are free; order-sensitive folds need a reversed twin, Hardening.)}
}
\end{recallbox}
